\thispagestyle{empty}
\vspace*{\fill}

\noindent
\textbf{Anatomía forense de una intrusión}\\
\textit{Disección pericial de un ataque a servidores de producción y respuesta a incidentes en sistemas Unix/Linux}

\vspace{1.5em}
\noindent
\textbf{Autor:} Alexis Real Santillán\\
\textbf{Función técnica original:} Administrador de Sistemas de Cómputo\\
\textbf{Contacto:} \href{mailto:alexisreal@gmail.com}{alexisreal@gmail.com}

\vspace{1.5em}
\noindent
\textbf{Nota sobre el origen del texto:}\\
Esta obra constituye una adaptación editorial del informe técnico y pericial elaborado por el Administrador de Sistemas de Cómputo y entregado a las autoridades de la institución tras la detección, contención y análisis del incidente de seguridad ocurrido en el centro de cómputo.

\vspace{1.5em}
\noindent
\textbf{Aviso legal y licencia de publicación:}\\
Esta obra se distribuye bajo una licencia internacional \textbf{Creative Commons Atribución-CompartirIgual 4.0 (CC BY-SA 4.0)}. Usted es libre de compartir (copiar y redistribuir el material en cualquier medio o formato) y adaptar (remezclar, transformar y construir a partir del material para cualquier propósito), bajo los términos de atribución adecuada y mantenimiento de la misma licencia.

\vspace{1.5em}
\noindent
\textbf{Nota editorial sobre anonimización y privacidad:}\\
El presente texto documenta una investigación pericial sobre incidentes ocurridos en sistemas reales de producción. Con el fin de permitir su divulgación pública para consulta técnica y profesional, se han aplicado estrictas directrices de anonimización:
\begin{itemize}
    \item Las direcciones IP del entorno de red han sido mapeadas dentro del bloque reservado para documentación \textbf{RFC 5737} (\texttt{192.0.2.0/24}, TEST-NET-1) y \textbf{RFC 2544}, preservando rigurosamente las relaciones de subred, las rutas y la correlación pericial de las sesiones.
    \item Los nombres de dominio y equipos han sido normalizados bajo la convención técnica (\texttt{ifm.lab.edu.mx}), asignando a los servidores identificadores clásicos inspirados en la nomenclatura habitual en centros de cómputo de la época (\texttt{helios}, \texttt{atlas}, \texttt{io}), de conformidad con el \textbf{RFC 2606}.
    \item Los nombres y cuentas de usuarios ajenos a la investigación han sido seudonimizados mediante identificadores representativos para resguardar la confidencialidad de datos personales.
\end{itemize}

\vspace{1.5em}
\noindent
\textit{Compuesto tipográficamente en \LaTeX{} utilizando fuentes TeX Gyre Pagella, TeX Gyre Heros e Inconsolata.}
\newpage
